\chapter{Linking, Loading and Relocation}

\section{Loaders}

So far every program was assumed to start at address 0. In reality the operating system
runs many programs and chooses where each one goes. The \textbf{loader} is the part of the OS
that:
\begin{enumerate}
  \item finds free memory at some address $\alpha$;
  \item copies the program's words into $\mathrm{MEM}[\alpha], \mathrm{MEM}[\alpha+4], \dots$;
  \item sets up registers (\code{\$30}, \code{\$31}, any inputs) and jumps to $\alpha$.
\end{enumerate}

\begin{lstlisting}[style=c]
// naive loader
alpha = findFreeRAM(n);
for (i = 0; i < n; i++)
    MEM[alpha + 4*i] = file[i];
$30 = alpha + n*4 + stackSize;   // or end of memory
$31 = returnToOS;
PC = alpha;
\end{lstlisting}

The OS itself is a program that is already loaded; conceptually it runs in a loop:
read a program name, load it, \code{jalr} to it, and continue when it returns.

\section{Relocation}

\paragraph{The problem.} Machine code assembled for address 0 is wrong at address $\alpha$
wherever it contains an \emph{absolute address}. In our subset that means exactly the
\code{.word label} lines: \code{.word foo} contains the address of \code{foo} assuming the
program starts at 0. Branch offsets are PC-relative, so they are fine.

\paragraph{Solution.} The loader adds $\alpha$ to every word that holds an address. But a
word in a binary file looks just like any other word, so the assembler must tell the loader
\emph{which} words to adjust. That information is stored in an \textbf{object file} format.
In CS 241 we use \textbf{MERL} (MIPS Executable Relocatable Linkable).

\section{The MERL format}

A MERL file consists of a 3-word header, the code, and a table of entries:

\begin{center}
\begin{tabular}{lll}
\toprule
offset & contents & meaning \\
\midrule
\code{0x00} & \code{0x10000002} & cookie: the instruction \code{beq \$0, \$0, 2} \\
\code{0x04} & \textit{endModule} & total length of the file in bytes \\
\code{0x08} & \textit{endCode} & length of header + code in bytes \\
\code{0x0c} & code \dots & the assembled program \\
\textit{endCode} & table \dots & relocation and external symbol entries \\
\bottomrule
\end{tabular}
\end{center}

The cookie is clever: it is a valid instruction that skips the next two words, so a MERL file
loaded at address 0 can be run directly, jumping over \textit{endModule} and \textit{endCode}.

Table entries:
\begin{center}
\begin{tabular}{llp{7cm}}
\toprule
type & format & meaning \\
\midrule
REL & \code{0x01}, \textit{loc} & word at \textit{loc} holds a relocatable address \\
ESR & \code{0x11}, \textit{loc}, \textit{len}, $c_1,\dots,c_{len}$ & external symbol
reference: word at \textit{loc} should hold the address of the imported symbol named
$c_1\dots c_{len}$ (one ASCII char per word) \\
ESD & \code{0x05}, \textit{loc}, \textit{len}, $c_1,\dots,c_{len}$ & external symbol
definition: this module exports a symbol with that name at address \textit{loc} \\
\bottomrule
\end{tabular}
\end{center}

All locations are relative to the beginning of the MERL file \emph{including} the 12-byte
header, i.e.\ the code is assembled as though it starts at address \code{0x0c}.

\begin{example}
Source:
\begin{lstlisting}[style=mips]
    lis $3
    .word ten
    jr $31
ten: .word 10
\end{lstlisting}
MERL output (each line one word):
\begin{lstlisting}[style=mips]
beq $0, $0, 2      ; 0x00 cookie
.word 0x24         ; 0x04 endModule
.word 0x1c         ; 0x08 endCode
lis $3             ; 0x0c
.word 0x18         ; 0x10 address of ten (0x0c + 0x0c)
jr $31             ; 0x14
.word 10           ; 0x18 ten
.word 0x01         ; 0x1c REL entry
.word 0x10         ; 0x20 location to relocate
\end{lstlisting}
The same file as a picture: the REL entry in the table points at the one word of code that
holds an address.
\begin{center}
\begin{tikzpicture}[>=Stealth,cell/.style={draw,minimum width=3.2cm,minimum height=0.5cm,font=\ttfamily\small,anchor=west},
  ad/.style={font=\ttfamily\scriptsize,anchor=east}]
  \node[cell,fill=gray!10] (r0) at (0,0) {0x10000002};
  \node[cell,fill=gray!10] (r1) at (0,-0.5) {0x24};
  \node[cell,fill=gray!10] (r2) at (0,-1.0) {0x1c};
  \node[cell] (r3) at (0,-1.5) {lis \$3};
  \node[cell,fill=blue!8] (r4) at (0,-2.0) {.word 0x18};
  \node[cell] (r5) at (0,-2.5) {jr \$31};
  \node[cell] (r6) at (0,-3.0) {.word 10};
  \node[cell,fill=yellow!12] (r7) at (0,-3.5) {0x01};
  \node[cell,fill=yellow!12] (r8) at (0,-4.0) {0x10};
  \foreach \k/\ad in {0/00,1/04,2/08,3/0c,4/10,5/14,6/18,7/1c,8/20} \node[ad] at (-0.1,-0.5*\k) {0x\ad};
  \draw[decorate,decoration={brace,amplitude=5pt}] (-1.0,-1.25) -- (-1.0,0.25) node[midway,left=6pt,font=\small]{header};
  \draw[decorate,decoration={brace,amplitude=5pt}] (-1.0,-3.25) -- (-1.0,-1.25) node[midway,left=6pt,font=\small]{code};
  \draw[decorate,decoration={brace,amplitude=5pt}] (-1.0,-4.25) -- (-1.0,-3.25) node[midway,left=6pt,font=\small]{table};
  \draw[->,thick,blue!60!black] (r8.east) -- ++(0.6,0) |- node[pos=0.25,right,font=\scriptsize]{REL: relocate the word at \code{0x10}} (r4.east);
  \draw[->] (r1.east) -- ++(0.4,0) node[right,font=\scriptsize]{endModule: total bytes};
  \draw[->] (r2.east) -- ++(0.4,0) node[right,font=\scriptsize]{endCode: start of table};
\end{tikzpicture}
\end{center}
\end{example}

\section{A relocating loader}

\begin{lstlisting}[style=c]
read cookie, endMod, endCode;
codeLen = endCode - 12;
alpha = findFreeRAM(codeLen + stack);
for (i = 0; i < codeLen; i += 4)
    MEM[alpha + i] = readWord();
i = endCode;
while (i < endMod) {
    format = readWord();
    if (format == 1) {               // REL
        rel = readWord();
        MEM[alpha + rel - 12] += alpha - 12;
        i += 8;
    } else error("unknown format");
}
\end{lstlisting}

Why $\alpha - 12$? The stored address already includes the header offset of 12 (code was
assembled starting at \code{0x0c}), but the loader places the first code word (not the header)
at $\alpha$. So the true address is $\mathit{stored} - 12 + \alpha$.

\begin{center}
\begin{tikzpicture}[>=Stealth,cell/.style={draw,minimum width=3.0cm,minimum height=0.55cm,font=\ttfamily\small},
  ad/.style={font=\scriptsize,anchor=east}]
  \node[font=\small] at (0,0.65) {in the MERL file};
  \node[cell] (f0) at (0,0) {lis \$3};
  \node[cell,below=0pt of f0,fill=blue!8] (f1) {.word 0x18};
  \node[cell,below=0pt of f1] (f2) {jr \$31};
  \node[cell,below=0pt of f2] (f3) {.word 10};
  \foreach \c/\ad in {f0/0x0c,f1/0x10,f2/0x14,f3/0x18} \node[ad,font=\ttfamily\scriptsize] at ([xshift=-2pt]\c.west) {\ad};
  \node[font=\small] at (6,0.65) {in memory};
  \node[cell] (m0) at (6,0) {lis \$3};
  \node[cell,below=0pt of m0,fill=blue!8] (m1) {.word $\alpha$+0x0c};
  \node[cell,below=0pt of m1] (m2) {jr \$31};
  \node[cell,below=0pt of m2] (m3) {.word 10};
  \foreach \c/\ad in {m0/$\alpha$,m1/$\alpha+4$,m2/$\alpha+8$,m3/$\alpha+12$} \node[ad] at ([xshift=-2pt]\c.west) {\ad};
  \foreach \i in {0,2,3} \draw[->,gray] (f\i.east) -- ([xshift=-1.1cm]m\i.west);
  \draw[->,thick,blue!60!black] (f1.east) -- node[above,font=\scriptsize]{$+(\alpha-12)$} ([xshift=-1.1cm]m1.west);
  \draw[->,dashed] (m1.east) to[out=0,in=0,looseness=1.6] node[right,font=\scriptsize]{\code{ten}} (m3.east);
\end{tikzpicture}

\small Loading at $\alpha$: ordinary words are copied; the REL word also gets $\alpha-12$ added.
\end{center}

\begin{important}[Position-independent code]
Code containing no absolute addresses (only PC-relative branches) is
\textbf{position-independent} and needs no relocation. Hidden absolute addresses that the
assembler cannot see (e.g.\ computing an address by arithmetic on a constant) will break
under relocation, so always produce addresses with \code{.word label}.
\end{important}

\section{Linking}

Large programs are split across files, and libraries (like the \code{print} routine or
\code{alloc.merl} used by our compiler) are assembled separately. A \textbf{linker} combines
object files into one, resolving references between them.

\paragraph{Directives.}
\begin{itemize}
  \item \code{.import name} declares that \code{name} is defined in another module.
        Using it in \code{.word name} generates an ESR entry (with value 0 in the code).
  \item \code{.export name} makes a label defined in this module visible to others
        (generates an ESD entry).
\end{itemize}
An assembler for MERL must check: an imported name must not also be a local label;
an exported name must be defined; a label used must be either defined or imported.

\subsection{Linking algorithm}
To link $m_1$ and $m_2$ (producing a module in which $m_2$'s code follows $m_1$'s):
\begin{enumerate}
  \item Check no symbol is exported by both (duplicate definition $\Rightarrow$ error).
  \item Let $\alpha = \mathit{endCode}(m_1) - 12$. Relocate $m_2$: add $\alpha$ to every
        location in $m_2$'s table, and add $\alpha$ to each word in $m_2$'s code marked REL.
  \item For each ESR in $m_1$ whose name is ESD-exported by $m_2$: write the ESD's address
        into the referenced word and change the ESR into a REL (it is now an internal
        address, still needing relocation at load time). Symmetrically for ESRs in $m_2$
        resolved by $m_1$'s ESDs.
  \item Concatenate: header, code of $m_1$, code of $m_2$, combined tables. Recompute
        \textit{endCode} and \textit{endModule}.
\end{enumerate}
ESRs that remain unresolved stay as ESRs, so linking can happen in stages; a final
executable should have none left.

\begin{center}
\begin{tikzpicture}[>=Stealth,font=\small,
  hd/.style={draw,fill=gray!10,minimum width=3cm,minimum height=0.5cm},
  cd/.style={draw,minimum width=3cm,align=center},
  tb/.style={draw,fill=yellow!12,minimum width=3cm,minimum height=0.6cm,align=center,font=\scriptsize}]
  % m1
  \node at (0,0.55) {$m_1$};
  \node[hd] (h1) at (0,0) {header};
  \node[cd,minimum height=1.6cm,fill=blue!5,below=0pt of h1] (c1) {code of $m_1$\\ \code{.word foo}};
  \node[tb,below=0pt of c1] (t1) {ESR \code{foo}};
  % m2
  \node at (0,-2.85) {$m_2$};
  \node[hd] (h2) at (0,-3.4) {header};
  \node[cd,minimum height=1.2cm,fill=green!6,below=0pt of h2] (c2) {code of $m_2$\\ \code{foo:} \dots};
  \node[tb,below=0pt of c2] (t2) {ESD \code{foo}, REL \dots};
  % linked
  \node at (6.5,0.55) {linked module};
  \node[hd] (h) at (6.5,0) {header (recomputed)};
  \node[cd,minimum height=1.6cm,fill=blue!5,below=0pt of h] (lc1) {code of $m_1$\\ \code{.word foo}};
  \node[cd,minimum height=1.2cm,fill=green!6,below=0pt of lc1] (lc2) {code of $m_2$\\ \code{foo:} \dots};
  \node[tb,minimum height=0.9cm,below=0pt of lc2] (lt) {REL (was ESR \code{foo}),\\ ESD \code{foo}$+\alpha$, $m_2$'s RELs $+\alpha$};
  \draw[->,gray] (c1.east) -- (lc1.west);
  \draw[->,gray] (c2.east) -- (lc2.west);
  \node[below=2pt of lt,font=\scriptsize,align=center] {$m_2$'s code now starts $\alpha=\mathit{endCode}(m_1)-12$\\ bytes later, so its addresses are shifted by $\alpha$};
  \draw[->,thick,red!60!black] ([yshift=-0.35cm]lc1.east) to[out=0,in=0,looseness=1.8] node[right,font=\scriptsize]{resolved} ([yshift=-0.2cm]lc2.east);
\end{tikzpicture}
\end{center}

\begin{lstlisting}[style=c]
alpha = m1.endCode - 12;
for (e : m2.table) e.loc += alpha;
for (e : m2.table) if (e.type == REL) m2.code[e.loc] += alpha;
for (d : m2.ESDs) d.addr += alpha;   // (same as loc for ESD)
for (r : m1.ESRs) if (m2 exports r.name) {
    m1.code[r.loc] = addressOf(r.name); r.type = REL; }
for (r : m2.ESRs) if (m1 exports r.name) {
    m2.code[r.loc] = addressOf(r.name); r.type = REL; }
\end{lstlisting}

\section{Static vs.\ dynamic linking}
\begin{aside}
In the course we do \emph{static} linking: all modules are merged into one file before
running. Real systems also use \emph{dynamic} linking (shared libraries, \code{.so}/\code{.dll}),
where references are resolved at load time or on first call, so library code is shared among
processes and can be updated without relinking. Real object formats (ELF, COFF, Mach-O) carry
the same kinds of information as MERL: code, relocation records, and symbol tables.
\end{aside}

\section{Summary}
\begin{itemize}
  \item The loader picks an address $\alpha$, copies code there, and must relocate absolute addresses.
  \item In our subset, only \code{.word label} produces absolute addresses; branches are relative.
  \item MERL = header (cookie \code{0x10000002}, endModule, endCode) + code + table of REL,
        ESR, ESD entries; locations include the 12-byte header.
  \item Relocation at load: $\mathrm{MEM}[\alpha+\mathit{loc}-12] \mathrel{+}= \alpha-12$.
  \item Linking: offset the second module, resolve ESRs against ESDs (turning them into RELs),
        concatenate, fix the header; duplicate exports are an error.
\end{itemize}
